\section{Related work and the scope of this comparison}
\label{sec:related-expanded}

\subsection{Combination, weight estimation and stacking}
Bates and Granger~\cite{bates} used past forecast errors to allocate weights to two competing forecasts. Their example establishes an important distinction for the present study: a combination can benefit from complementary errors without either component being universally best. It does not establish that estimating more flexible weights is beneficial with a small memory. Our shared-expert comparisons therefore ask whether a particular allocation procedure improves prediction, rather than treating the presence of multiple experts as evidence for contextual learning.

Clemen's review~\cite{clemen} documents both the practical interest in combining forecasts and the robustness of simple combinations. This motivates keeping EqualMean visible even when it is not the primary mechanism-isolating control. An estimated procedure should not be judged only against a deliberately weak single expert. Conversely, a gain over equal weights cannot identify why the allocation improved: recency, feature matching, temperature, regularization and historical selection are separate changes in the present implementation.

Timmermann~\cite{timmermann} organizes the combination problem around error dependence, instability and weight-estimation uncertainty. This provides a more useful explanation for our design than a generic claim that ensembles reduce variance. All five experts receive the same observations, so their errors can be strongly related. The study does not assume independent experts, nor does the use of five columns create five independent replications. The Gram matrix used by stacking retains cross-products rather than replacing dependence by a diagonal approximation.

Smith and Wallis~\cite{smithwallis} analyze finite-sample estimation error as an explanation of the forecast-combination puzzle. In our setting, a maximum of twelve remembered origins and five expert coefficients makes that concern concrete. A flexible allocation can fit an unstable relationship among historical errors. Our uniform floor and positive regularization restrict that flexibility; neither device proves that its estimation cost is smaller than its predictive benefit. The observed unfavorable fixed-context comparison is consistent with that concern but does not identify its cause.

Wolpert's stacked generalization~\cite{wolpert} learns from predictions made on observations not used to train the corresponding first-stage predictors. This information separation is central to the current comparison. Predictions used to assess a weight rule must have been issued before their targets became available. Our temporal blocks implement that separation for the stated horizon. They do not reproduce arbitrary shuffled folds, and the origin schedule remains part of the estimator rather than an interchangeable evaluation detail.

Breiman's stacked regressions~\cite{breiman} use cross-validation predictions and nonnegative regression coefficients to construct combinations. This is direct prior art for learning constrained linear combinations from held-out errors. Our stacking candidate additionally restricts coefficients to sum to one and penalizes distance from equal weights. Its historical absolute-error selection criterion also differs from its squared-error fitting objective. These differences define an implementation to evaluate; they do not support a priority claim for nonnegative stacking itself.

Hansen's least-squares model averaging~\cite{hansen} selects weights using a Mallows criterion and proves asymptotic properties for its specified model class and assumptions. That theoretical result does not transfer to our short-memory estimator by analogy. We use a different criterion, fitted experts, a capped origin schedule and dependent temporal records. The present theoretical statements are instead finite algebraic properties: uniqueness of the regularized quadratic solution, temporal measurability, convex quantile ordering and a limited absolute-error inequality.

Raftery, K\'arn\'y and Ettler~\cite{dma} develop dynamic model averaging with evolving model parameters and model uncertainty. Their work makes clear that time-varying allocation is established research. Our procedure has no equivalent latent-state probability model and does not claim to approximate that algorithm. We retain a small set of realized expert errors and recompute weights at the requested prefix. Replaying those fits is distinct from a native recursive implementation of dynamic model averaging.

\subsection{Features, conditional performance and current alternatives}
FFORMA~\cite{fforma} learns a feature-dependent mapping to forecast weights across reference series. Arbitrated forecasting ensembles~\cite{ade} estimate expert competence and use it to form adaptive predictions. Both are more relevant conceptual comparators than unrelated large models because they address the allocation mechanism. Our recurrence rule estimates relevance from completed origins within one series rather than fitting a global competence model. A fair empirical comparison would need the original training resources, tuning rules, expert roster and information boundaries of each alternative; we have not executed it here.

The broad review by Wang and colleagues~\cite{wangreview} places forecast combination within a literature spanning statistical weighting, optimization and learning. The current study occupies a narrow part of that landscape: small local expert sets, a finite memory and complete historical target blocks. It does not cover hierarchical reconciliation, cross-series representation learning or all combinations of probabilistic distributions. Describing the procedure as better than all published forecast-combination methods would therefore overstate both the comparator roster and the empirical population.

A recent alternative is corrected forecast combinations by Liu and Vasnev~\cite{corrected}, whose September 2026 revision considers predictable serial dependence in combined forecast errors. Their correction uses an observable previous error rather than treating all remaining combined error as innovation. This offers a concrete candidate for another experiment, particularly when a completed residual remains relevant at the next query. It has not been implemented or scored in this study. Its results cannot be claimed as our results, and our outcomes do not establish superiority over it.

Giacomini and White~\cite{giacomini} distinguish conditional predictive ability from unconditional average performance. This distinction matters because our hypothesis concerns context, while our headline outcome is an average over admitted series. A lower aggregate MASE does not establish that the four context features identify which expert will perform best next. The reported experiment supplies matching controls and selection records, but no conditional predictive-ability test or causal identification of the feature mechanism. Such a test would require an explicit design and suitable forecast records.

\subsection{Temporal evaluation and empirical populations}
Tashman~\cite{tashman} treats splitting rules, forecast origins, rolling windows and coefficient updates as explicit choices in forecasting evaluation. Our source reconstruction, prefix cap and nonoverlapping outer target blocks follow that need for a fully specified experiment. Three eligible origins provide more than one temporal location for many series, but remain a limited view of a potentially long process. They do not demonstrate behavior under every change in frequency, update policy, horizon or seasonal lag.

Bergmeir and Ben\'itez~\cite{berg2012} examine how time-series prediction and overlapping lagged observations complicate evaluation. Their work motivates stating the available information rather than applying a generic cross-validation label. In our construction, an inner training block is admitted only after its last target arrives. This removes direct future-target use at the query. It does not remove dependence between neighboring observations, overlapping fitting prefixes or related series, and it is not presented as proof that all uncertainty calculations become independent.

Bergmeir, Hyndman and Koo~\cite{berg2018} show that ordinary cross-validation can be valid for autoregressive prediction under specified error conditions. We therefore do not claim that shuffled cross-validation is universally invalid for time series. Our chronological procedure addresses a different practical contract: every candidate is evaluated with information that would have been observable before its own validation target. That contract remains useful even when a separate theorem would justify another design for another model class.

The Monash archive~\cite{monash} makes many forecasting collections available in a common format. We identify its weekly and monthly releases by exact source bytes, individual sequence digests and eligibility decisions. This supports checking which series were actually included. It does not turn an archive into a random sample of all operational forecasting tasks. Shared subject matter, related source systems and exclusions can still make an exchangeable-series approximation optimistic.

The M3 study~\cite{m3study} and M4 competition~\cite{m4study} concern their own competition protocols and populations. Our re-use of monthly and weekly sequences changes the available prefix, origin schedule, comparator roster and aggregation rule. Thus, a method with lower MASE in our table has not earned an official competition placement. We report the two collections separately, disclose exclusions and retain AutoTheta and AutoETS so that the proposed combinations face automatic statistical alternatives as well as matching ablations.

\subsection{Scores, uncertainty and claims}
Gneiting and Raftery~\cite{proper} explain why probabilistic evaluation should reward useful concentration while penalizing failures to contain outcomes. Their interval-scoring perspective motivates retaining MSIS beside empirical coverage. A wide Gaussian interval can cover most targets and still provide a poor interval score. In the present comparison, the combination methods' point gains and interval losses are therefore jointly material. Coverage alone is not treated as a quality ranking or a calibration certificate.

Dem\v{s}ar~\cite{demsar} discusses comparisons of learning algorithms over multiple datasets, where the sampling unit differs from individual within-dataset observations. That distinction motivates our series-level summaries, but the paper's classifier-comparison recommendations are not automatically a theorem for this forecasting design. We do not manufacture a larger sample by counting every lead and every remembered fit independently. Our bootstrap preserves all admitted origins of a resampled series and remains conditional on a cross-series exchangeability approximation.

Diebold's discussion of predictive-accuracy tests~\cite{dmreview} emphasizes the need to match a comparison to the forecasting question. Our paired-series intervals are not Diebold--Mariano tests on a long stationary loss-difference sequence. They are a descriptive resampling analysis of series summaries after the outcomes were available. Keeping that distinction visible avoids attaching a familiar test name to a different estimand or interpreting an interval that excludes zero as prospective population confirmation.

These relationships define the contribution claimed here. The implementation combines established ingredients in an auditable finite-memory experiment, makes its nested information order explicit, and supplies complete outcomes for specified collections. The favorable results concern temporal selection relative to its recency control. The unfavorable fixed-context, monthly AutoTheta and interval comparisons remain part of the same claim. The literature supplies better questions and stronger candidate comparisons; it does not authorize a broader conclusion than the experiment actually tests.
